\documentclass[10pt,a4paper]{article}

% ---------- Packages ----------
\usepackage[utf8]{inputenc}
\usepackage[T1]{fontenc}
\usepackage[english]{babel}

\usepackage[
    a4paper,
    margin=1.8cm,
    top=1.2cm,
    bottom=1.5cm,
    includehead,
    headheight=1.9cm,
    headsep=0.45cm
]{geometry}

\usepackage{graphicx}
\usepackage[table]{xcolor}
\usepackage{titlesec}
\usepackage{enumitem}
\usepackage{fancyhdr}
\usepackage{array}
\usepackage{tabularx}
\usepackage{xurl}
\usepackage{hyperref}

% ---------- Colors ----------
\definecolor{mainblue}{RGB}{25,55,110}
\definecolor{lightgray}{RGB}{245,245,245}

% ---------- Links ----------
\hypersetup{
    colorlinks=true,
    linkcolor=mainblue,
    urlcolor=mainblue,
    citecolor=mainblue
}

% ---------- Header ----------
\pagestyle{fancy}
\fancyhf{}
\renewcommand{\headrulewidth}{0.4pt}
\renewcommand{\footrulewidth}{0pt}

\fancyhead[L]{%
    \begin{minipage}[c]{0.30\textwidth}
        \raggedright
        \IfFileExists{uzi_logo.png}{\includegraphics[width=3.2cm]{\detokenize{uzi_logo.png}}}{\textbf{UZH}}
    \end{minipage}
}

\fancyhead[C]{%
    \begin{minipage}[c]{0.36\textwidth}
        \centering
        \textbf{\small Thesis Proposal}\\[-0.1em]
        \scriptsize Department of Informatics
    \end{minipage}
}

\fancyhead[R]{%
    \begin{minipage}[c]{0.30\textwidth}
        \raggedleft
        \IfFileExists{everlogo.png}{\includegraphics[width=3.0cm]{\detokenize{everlogo.png}}}{\textbf{Everllence}}
    \end{minipage}
}

\fancyfoot[C]{\thepage}

% ---------- Section Formatting ----------
\titleformat{\section}
  {\color{mainblue}\normalfont\large\bfseries}
  {}
  {0em}
  {}

\titlespacing*{\section}{0pt}{0.6em}{0.25em}

% ---------- Compact Lists ----------
\setlist[itemize]{noitemsep, topsep=2pt, leftmargin=1.2em}
\setlist[enumerate]{noitemsep, topsep=2pt, leftmargin=1.4em}

% ---------- Tables ----------
\renewcommand{\arraystretch}{1.08}
\setlength{\tabcolsep}{6pt}

\begin{document}

% ---------- Title ----------
\begin{center}
    {\Large \textbf{Retrieve and Verify: Agentic Retrieval-Augmented Generation for\\[0.15em]Clause-Level Compliance Checking against Public Engineering Standards}}\\[0.3em]
    {\normalsize Industry Collaboration with Everllence SE (formerly MAN Energy Solutions)}
\end{center}

\vspace{0.35em}

% ---------- Student Information ----------
\noindent
\begin{tabularx}{\textwidth}{>{\bfseries}l X >{\bfseries}l X}
Student Name: & Necati Furkan Colak       & Student ID:    & 24746323 \\
Program:      & MSc Computer Science      & Semester:      & Spring 2026 \\
Supervisor:   & [Supervisor Name]         & Co-Supervisor: & [Everllence Contact] \\
Email:        & necatifurkan.colak@uzh.ch & Date:          & \today \\
\end{tabularx}

\vspace{0.5em}
\hrule
\vspace{0.5em}

\section{Abstract}

Engineering organizations must check that procedures, design descriptions, and work instructions follow the standards that apply to them. Today this check is manual: a reviewer reads the document and the standard side by side, which is slow and misses clauses. Retrieval-augmented generation (RAG) assistants answer questions about standards, but compliance checking is a different task. The system has to find the clauses that apply to each document section on its own, compare two texts, and assign a label (satisfied, violated, or insufficient evidence) with citations a reviewer can inspect. Earlier research automated this with rule-based pipelines that extract requirements from regulatory text \cite{zhang2016}; building the rules takes expert effort for every standard, and free-text procedures fit the rule format poorly. Agentic retrieval-augmented generation removes the rule-engineering step \cite{yao2023react, singh2025agentic}, and its accuracy on this task has not been measured. This thesis builds an agent that segments a document, retrieves candidate clauses per section, classifies compliance per clause, verifies its own labels, and writes a cited report. Evaluation uses a public standards corpus and expert-labeled cases, and measures clause retrieval recall, label accuracy, citation correctness, and reviewer time saved. All data is public, so the method transfers to any organization that reviews documents against standards.

\section{Motivation and Problem Statement}

\textbf{Operational context.} Standards govern engineering work: design codes, safety norms, and quality procedures. Before a document is released, someone must confirm that it follows the applicable standards. Reviewers do this by reading both texts, which is repetitive and depends on the reviewer remembering which clauses apply. Missed clauses surface later as audit findings or rework.

\textbf{Why question answering does not solve this.} RAG assistants answer the questions a user thinks to ask. Compliance review inverts the direction: the system must decide which clauses apply to each part of a document without being asked, then judge agreement between two texts. This requires per-section retrieval, an explicit comparison step, and an ``insufficient evidence'' outcome, because a wrong ``compliant'' label is worse than no label. Rule-based compliance checking research addressed the same task by extracting requirements into formal rules and matching them against structured design data \cite{zhang2016}. This works where rules exist, and it is costly to extend, because every new standard and every free-text procedure needs new rules. LLM agents supply the missing control flow of retrieval, comparison, and self-checking without hand-built rules \cite{yao2023react, singh2025agentic}.

\textbf{The open question.} No study measures how well an agentic RAG pipeline performs full clause-level compliance checking: coverage of the applicable clauses, label quality per clause, and reliability of the citations. This thesis provides those measurements on public data, so the design stays reproducible and independent of any single company.

\section{Research Questions}

\begin{enumerate}
    \item \textbf{RQ1 (primary):} How accurately does an agentic RAG pipeline identify the standard clauses applicable to a given document section and classify each as satisfied, violated, or insufficient evidence, compared with expert judgments?
    \item \textbf{RQ2:} Which pipeline components (per-section query generation, clause reranking, self-verification of labels) contribute most to accuracy?
    \item \textbf{RQ3:} How reliable are the produced citations: does each cited clause justify the assigned label?
    \item \textbf{RQ4:} How much reviewer time does the assistant save in a small expert study, and which error types remain?
\end{enumerate}

\section{Methodology}

\begin{itemize}
    \item \textbf{Literature review:} RAG foundations \cite{lewis2020rag}; agentic architectures \cite{yao2023react, singh2025agentic}; automated compliance checking \cite{zhang2016}; RAG evaluation \cite{es2024ragas, ragindustry2025}.
    \item \textbf{Data:} A public standards corpus (for example openly accessible Eurocode national annexes, public ISO and IEC extracts, or public safety regulations) plus 30 to 60 short engineering documents (procedures, design descriptions) taken from public sources or written for the benchmark. Experts label the applicable clauses and the per-clause compliance for each document section, giving a gold set of roughly 300 to 600 clause-document pairs with a development and a held-out test split. No proprietary data is needed.
    \item \textbf{System:} An agent that (1) segments the input document, (2) generates retrieval queries per section, (3) retrieves and reranks candidate clauses with hybrid BM25 and dense retrieval, (4) classifies each pair as satisfied, violated, or insufficient evidence, (5) verifies its own label against the quoted clause text, and (6) writes a cited report. Baselines: a single-pass RAG prompt without the per-section agent steps, and a closed-book LLM. Ablations remove self-verification and reranking (RQ2).
    \item \textbf{Models and tools:} One GPT-4-class API model (for example via Azure OpenAI) and one smaller open-weight model to check that findings hold across model sizes; LlamaIndex or LangGraph with a standard vector database.
    \item \textbf{Evaluation:} Clause retrieval recall and precision per section; per-clause label accuracy and F1 against expert labels; citation correctness; error rates of the ``insufficient evidence'' class in both directions; paired significance tests between system variants; a timed study with 3 to 5 engineers comparing manual review with assistant-supported review on held-out documents.
\end{itemize}

\section{Expected Contribution}

\begin{enumerate}
    \item \textbf{For research:} Measured performance of agentic RAG on clause-level compliance checking, a task that differs from question answering in retrieval direction and output structure. The work replaces the hand-built rules of earlier compliance checking systems \cite{zhang2016} with retrieval and generation, and it quantifies what this trade costs in accuracy.
    \item \textbf{For practice:} A tested review-support method for compliance work in quality management, safety, and certification, applicable wherever documents must be checked against standards.
    \item \textbf{Open deliverables:} The labeled clause-document benchmark, prompts, and pipeline code. All inputs are public, so the results can be reproduced.
\end{enumerate}

\section{Proposed Timeline (6 months)}

\begin{tabularx}{\textwidth}{|p{2.6cm}|X|}
\hline
\rowcolor{lightgray}
\textbf{Period} & \textbf{Planned Work} \\
\hline
Month 1 & Literature review; standards corpus selection. \\
\hline
Month 2 & Document set preparation and expert labeling; baseline implementations. \\
\hline
Month 3 & Agent pipeline and the component ablations. \\
\hline
Month 4 & Full experimental runs; statistical analysis (RQ1 to RQ3). \\
\hline
Month 5 & Expert time study (RQ4); error analysis. \\
\hline
Month 6 & Thesis writing; code release; final presentation. \\
\hline
\end{tabularx}

\begingroup
\renewcommand{\refname}{\texorpdfstring{\color{mainblue}}{}Preliminary References}
\scriptsize
\begin{thebibliography}{9}
\setlength{\itemsep}{0pt}

\bibitem{zhang2016}
J.\ Zhang and N.\ El-Gohary.
\textit{Semantic NLP-Based Information Extraction from Construction Regulatory Documents for Automated Compliance Checking}.
Journal of Computing in Civil Engineering, 2016.

\bibitem{lewis2020rag}
P.\ Lewis et al.
\textit{Retrieval-Augmented Generation for Knowledge-Intensive NLP Tasks}.
NeurIPS, 2020.

\bibitem{yao2023react}
S.\ Yao et al.
\textit{ReAct: Synergizing Reasoning and Acting in Language Models}.
ICLR, 2023.

\bibitem{singh2025agentic}
A.\ Singh et al.
\textit{Agentic Retrieval-Augmented Generation: A Survey on Agentic RAG}.
arXiv:2501.09136, 2025.

\bibitem{es2024ragas}
S.\ Es et al.
\textit{RAGAS: Automated Evaluation of Retrieval Augmented Generation}.
EACL, 2024.

\bibitem{ragindustry2025}
\textit{Retrieval-Augmented Generation in Industry: An Interview Study on Use Cases, Requirements, Challenges, and Evaluation}.
arXiv:2508.14066, 2025.

\end{thebibliography}
\endgroup

\vspace{0.1em}

\noindent
\begin{tabularx}{\textwidth}{X X}
\textbf{Student Signature:} & \textbf{Supervisor Signature:} \\[0.9em]
\rule{0.42\textwidth}{0.4pt} & \rule{0.42\textwidth}{0.4pt} \\
\end{tabularx}

\end{document}
